\section{The Bailout Protocol}\label{sec:bailout}

The Bailout Protocol is the mandatory exception-escalation mechanism of the \bxthree{} Framework. It is not an error-handling routine selected at runtime, nor an optional escalation pathway available to a system that chooses to invoke it. It is an architectural compulsion: any unresolved exception state encountered by any node in a \bxthree{} system \emph{must} propagate upward through the accountability chain until it reaches a named human principal, and no machine actor along that chain may absorb, resolve, or suppress the exception. This section provides the complete formal specification: the deterministic state machine governing exception propagation, the precise conditions governing each transition, the bail-out trigger condition, the escalation algorithm, the bypass mechanism enforcing the no-machine-actor property, and the timeout handling that guarantees liveness under all failure conditions.

The protocol operates against the backdrop of the \bxthree{} three-layer model. The \purpose{} layer carries intent, judgment, and final authority over any decision exceeding a node's provisioned scope; it defines the operating range $R(n)$ for each node $n$ and receives all escalated exceptions. The \bounds{} layer carries bounded reasoning and constrained execution within $R(n)$; it evaluates whether proposed actions fall within the node's authority and initiates bailout when they do not. The \fact{} layer carries deterministic enforcement and forensic auditability; it verifies proposed actions against the authorization ledger, enforces the transition relation of the state machine, and maintains the append-only Forensic Ledger recording every protocol event as an immutable architectural record. Each spawned node carries an immutable parent pointer $\alpha(n)$, and every node carries a reference to its human accountability anchor $h(n)$ — a named human principal or institutional actor — established at node provisioning and never modified for the node's lifetime. The Bailout Protocol exploits this structure to guarantee that exception states propagate to $h(n)$ without passing through any intermediate machine actor as a terminus.

% --------------------------------------------------------------
\subsection{State Machine Formalism}\label{sec:state-machine}

The Bailout Protocol is formally modeled as a deterministic finite state machine
\begin{equation}
\mathcal{B} \;=\; \bigl(Q,\; q_0,\; \Sigma,\; \rightarrow,\; Q_F\bigr)
\end{equation}
where:

\begin{itemize}\itemsep0em
\item $Q = \{\,\text{IDLE},\; \text{BOUNDED},\; \text{BAIL\_PENDING},\; \text{ESCALATING},\; \text{HUMAN\_ACKNOWLEDGED},\; \text{RESOLVED}\,\}$ is the finite set of protocol states;
\item $q_0 = \text{IDLE}$ is the designated initial state, entered at node startup and after every successful directive resolution;
\item $\Sigma$ is the input alphabet, comprising the trigger event $e_{\text{trigger}}$, authorization confirmation $e_{\text{auth}}$, timeout signal $e_{\text{timeout}}$, human acknowledgment $e_{\text{ack}}$, binding resolution directive $e_{\text{resolve}}$, and rejection directive $e_{\text{reject}}$;
\item $\rightarrow\;\subseteq Q \times \Sigma \times Q$ is the deterministic transition relation (defined exhaustively in Table~\ref{tab:bailout-transition});
\item $Q_F = \{\text{RESOLVED}\}$ is the set of terminal accepting states.
\end{itemize}

A node's state $q \in Q$ is stored in the node's Fact Layer worksheet and is observable by its parent node and by the Bailout Monitor. State updates are atomic; no intermediate or partially-executing states are observable by any external actor. The machine is \emph{deterministic}: for every $(q, \sigma) \in Q \times \Sigma$, there exists at most one $q' \in Q$ such that $(q, \sigma, q') \in \rightarrow$. Determinism is enforced by the Bailout Monitor's priority function $\pi$ (defined in \S\ref{sec:bailout-trigger}), which totally orders simultaneous trigger conditions and selects the highest-priority enabled transition before committing any state change. If no transition is defined for a $(q, \sigma)$ pair, the machine stalls — this is a protocol violation and is recorded as a \texttt{stall} entry in the Forensic Ledger.

\begin{table}[htbp]
\caption{Bailout Protocol state transition relation $\rightarrow\;\subseteq Q \times \Sigma \times Q$. Rows are current state; columns are input events. Entries denote the next state. Blank cells indicate that no transition is defined for that $(q, \sigma)$ pair — the machine stalls and the Bailout Monitor records a protocol violation.}
\label{tab:bailout-transition}
\begin{tabular}{@{}lcccccc@{}}
\toprule
 & \multicolumn{6}{c}{Input event} \\
\cmidrule(l){2-7}
Current state & $e_{\text{trigger}}$ & $e_{\text{auth}}$ & $e_{\text{timeout}}$ & $e_{\text{ack}}$ & $e_{\text{resolve}}$ & $e_{\text{reject}}$ \\
\midrule
IDLE          & BOUNDED             & —                 & —                   & —                & —                   & —                 \\
BOUNDED       & —                   & HUMAN\_ACKNOWLEDGED & BAIL\_PENDING       & —                & —                   & —                 \\
BAIL\_PENDING & —                   & —                 & ESCALATING          & —                & —                   & —                 \\
ESCALATING    & —                   & HUMAN\_ACKNOWLEDGED & ESCALATING         & IDLE             & RESOLVED            & RESOLVED          \\
HUMAN\_ACKNOWLEDGED & —           & —                 & —                   & IDLE             & RESOLVED            & IDLE              \\
RESOLVED      & —                   & —                 & —                   & —                & —                   & —                 \\
\bottomrule
\end{tabular}
\end{table}

\subparagraph{State definitions.}

\begin{itemize}\itemsep0em
\item[\textbf{IDLE.}] The node is operating within its provisioned authority. No exception condition is active. The Fact Layer processes normal action requests through the standard dispatch path. The node may transition to BOUNDED at any time upon receipt of a trigger event satisfying TC (defined in \S\ref{sec:bailout-trigger}).

\item[\textbf{BOUNDED.}] The node has detected an exception condition and is evaluating whether it is autonomously resolvable within current authority. The Bounds Engine is actively executing its resolution procedure. All machine-actor resolution paths remain open. If the Bounds Engine produces a resolution $a$ that the Fact Layer approves before $T_{\text{bounds}}$ expires, the node returns to IDLE via T2. If all resolution paths are exhausted or $T_{\text{bounds}}$ expires, the state advances to BAIL\_PENDING via T3.

\item[\textbf{BAIL\_PENDING.}] The node has determined — either through explicit \texttt{bailout\_signal} from the Bounds Engine or through $T_{\text{bounds}}$ expiration without a \fact-layer-approved resolution — that the exception exceeds its current authority. The Bailout Protocol is formally activated, but propagation to $h(n)$ has not yet been initiated. This is the last architectural point at which a machine actor could theoretically suppress escalation. It cannot: entry to BAIL\_PENDING triggers an immediate atomic transition to ESCALATING via T4, with no dwell time and no machine-actor intervention permitted.

\item[\textbf{ESCALATING.}] The exception is propagating upward through the parent chain toward $h(n)$. Each intermediate node receives the exception, appends its diagnostic context to the propagation chain via \texttt{ForwardException} (detailed in Algorithm~\ref{alg:forward-exception}), and forwards the exception to its parent via \texttt{SendToParent} without attempting local resolution. The state remains ESCALATING until either (a) the human anchor acknowledges receipt ($e_{\text{ack}}$: T5 transitions to HUMAN\_ACKNOWLEDGED), (b) the human anchor issues a binding directive ($e_{\text{resolve}}$ or $e_{\text{reject}}$: T7 transitions to RESOLVED), or (c) all pathways time out after $N_{\max}$ retries, propagating a \texttt{parent\_unreachable} event directly to $h(n)$ via an alternate functional pathway.

\item[\textbf{HUMAN\_ACKNOWLEDGED.}] The human anchor has received the exception, observed the diagnostic context, and acknowledged it. The node awaits a directive. No machine actor may issue, simulate, or substitute a directive in this state. The node remains in this state until the human principal issues \texttt{ACK} ($e_{\text{ack}}$: T8 returns to IDLE) or a binding directive $e_{\text{resolve}}\) ($e_{\text{reject}}$: T9 enters RESOLVED).

\item[\textbf{RESOLVED.}] The human principal has issued a binding directive and the protocol has recorded it in the authorization ledger. The directive is one of: \texttt{ACK} (continue with current parameters, return to IDLE), \texttt{RESOLVE} (execute a specific binding resolution, enter quiescent post-execution state), or \texttt{REJECT} (disallow the proposed action, enter a quiescent error state). This state is terminal for the current protocol run.
\end{itemize}

\subparagraph{State invariants.} The following invariants hold for all reachable states of $\mathcal{B}$:

\begin{align}
\text{INV}_1 &: \forall n \in \text{nodes},\; \text{state}(n) \in Q \setminus \{\text{ESCALATING}\} \iff \neg\text{active\_bailout}(n). \tag{INV-1}
\end{align}
That is, a node occupies a non-escalating state \emph{iff} it has no active bailout propagation in flight. A node may carry a bailout history (a ledger entry recording a past propagation) while in IDLE, but no active propagation may be underway.

\begin{align}
\text{INV}_2 &: \forall n \in \text{nodes},\; \text{state}(n) = \text{ESCALATING} \implies \text{pathway}(n) \neq \emptyset. \tag{INV-2}
\end{align}
Every node in ESCALATING state has at least one functional communication pathway to its parent, selected by the Pathway Health Registry before the first \texttt{SendToParent} call.

\begin{align}
\text{INV}_3 &: \forall n \in \text{nodes},\; \text{parent}(n) = \alpha(n) \implies h(\alpha(n)) = h(n). \tag{INV-3}
\end{align}
The human anchor reference is transitively conserved up the parent chain: every node shares the same human anchor $h(n)$, established at the root node and propagated to all descendants at spawn time.

\subparagraph{Transition definitions.} The six enabled transitions of $\mathcal{B}$ are:

\begin{align}
\text{T1} &:\; \text{IDLE} \xrightarrow{e_{\text{trigger}}} \text{BOUNDED} \tag{T1} \\
\text{T2} &:\; \text{BOUNDED} \xrightarrow{e_{\text{auth}}} \text{HUMAN\_ACKNOWLEDGED} \tag{T2} \\
\text{T3} &:\; \text{BOUNDED} \xrightarrow{e_{\text{timeout}}} \text{BAIL\_PENDING} \tag{T3} \\
\text{T4} &:\; \text{BAIL\_PENDING} \xrightarrow{e_{\text{timeout}}} \text{ESCALATING} \tag{T4} \\
\text{T5} &:\; \text{ESCALATING} \xrightarrow{e_{\text{ack}}} \text{HUMAN\_ACKNOWLEDGED} \tag{T5} \\
\text{T6} &:\; \text{ESCALATING} \xrightarrow{e_{\text{timeout}}} \text{ESCALATING} \tag{T6} \\
\text{T7} &:\; \text{ESCALATING} \xrightarrow{e_{\text{resolve}} \lor e_{\text{reject}}} \text{RESOLVED} \tag{T7} \\
\text{T8} &:\; \text{HUMAN\_ACKNOWLEDGED} \xrightarrow{e_{\text{ack}}} \text{IDLE} \tag{T8} \\
\text{T9} &:\; \text{HUMAN\_ACKNOWLEDGED} \xrightarrow{e_{\text{resolve}} \lor e_{\text{reject}}} \text{RESOLVED} \tag{T9}
\end{align}

T6 is a self-loop on ESCALATING: each occurrence increments the retry counter and doubles $T_{\text{prop}}$ up to $T_{\text{cap}}$. After $N_{\max}$ consecutive T6 firings without $e_{\text{ack}}$, the escalation algorithm fires \texttt{parent\_unreachable} directly to $h(n)$ via an alternate pathway (Algorithm~\ref{alg:escalate}, line~\ref{line:parent-unreachable}).

% --------------------------------------------------------------
\subsection{Bail-Out Trigger Condition}\label{sec:bailout-trigger}

The bail-out trigger condition (TC) is the mechanism by which the \bxthree{} system detects when a node has encountered a condition outside its provisioned operating range that it cannot autonomously resolve. TC is the decisive condition that drives T1, causing a node to enter BOUNDED state. Rather than requiring system designers to enumerate specific error codes in advance — a practice that provides coverage only for anticipated failure modes and leaves unanticipated ones invisible — the Bailout Protocol defines the trigger in terms of \emph{observational boundary violation}: a node escalates precisely when its observable input space contains a condition outside its provisioned operating range $R(n)$ and the node cannot produce a \fact-layer-approved resolution within $T_{\text{bounds}}$.

\subparagraph{Primary trigger condition.} The trigger $\text{TRIGGER}(n, x)$ is defined as a conjunction of three sub-conditions, all of which must hold simultaneously:
\begin{equation}
\text{TRIGGER}(n, x) \;\equiv\; \underbrace{x \notin R(n)}_{(i)} \;\land\; \underbrace{\neg\text{resolved}(n, x)}_{(ii)} \;\land\; \underbrace{\text{confidence}(n, x) < \tau}_{(iii)} \tag{TC}
\end{equation}
where:

\begin{itemize}\itemsep0em
\item[(i)] $x$ is an element of the node's observable input space that falls outside the provisioned operating range $R(n)$. This is the \emph{bound signal} condition: the node observes something it was not provisioned to handle autonomously.
\item[(ii)] $x$ has not been previously resolved by the node within the current operational session. This prevents re-trigger on known-exception inputs that have already received human adjudication.
\item[(iii)] $\text{confidence}(n, x) < \tau$ is the node's self-assessed confidence that it can produce a correct resolution for $x$ without human input, where $\tau$ is the provisioned confidence threshold. The confidence function $\text{confidence} : \text{node} \times \text{input} \rightarrow [0, 1]$ is defined at node provisioning and is \emph{not} adjustable at runtime by any machine actor. This prevents a machine actor from inflating its own confidence to suppress a warranted escalation.
\end{itemize}

\subparagraph{Trigger priority function.} When multiple trigger conditions fire simultaneously for node $n$, they are totally ordered by the priority function $\pi(n, x)$, which ensures deterministic transition selection:
\begin{equation}
\pi(n, x) \;=\; w_1 \cdot \text{severity}(x) \;+\; w_2 \cdot \bigl(1 - \text{confidence}(n, x)\bigr) \;+\; w_3 \cdot \text{distance}\bigl(n, h(n)\bigr) \tag{PF}
\end{equation}
where $\text{severity}(x) \in [0, 1]$ is the assigned severity of exception $x$ (higher values indicate greater potential impact), $\text{distance}(n, h(n))$ is the number of hops from $n$ to the human anchor in the parent chain (longer distances increase priority because propagation latency is greater), and $w_1, w_2, w_3$ are provisioned constants satisfying $w_1 > w_2 > w_3 \geq 0$. The highest-priority trigger fires first. Ordering is enforced by the Bailout Monitor before any state transition is committed to the Forensic Ledger.

\subparagraph{Secondary trigger: explicit \texttt{bailout\_signal}.} In addition to TC, the Bounds Engine may emit an explicit \texttt{bailout\_signal} at any time, regardless of confidence scores. This mechanism handles the case where the Bounds Engine detects a condition — for example, a logical contradiction between two inputs within $R(n)$ — that it can recognize as unresolvable even though both inputs are individually within the operating range. The \texttt{bailout\_signal} bypasses TC entirely and immediately initiates T3:
\begin{equation}
\texttt{bailout\_signal}(x) \;\implies\; \text{T3 fires for node } n \text{ with exception } x. \tag{TC-secondary}
\end{equation}

% --------------------------------------------------------------
\subsection{Escalation Algorithm}\label{sec:escalation-algorithm}

When the state machine enters ESCALATING (T4), the protocol executes the \texttt{Escalate} algorithm to propagate the exception to $h(n)$ along the parent chain. The originating node's Bailout Monitor executes \texttt{Escalate}; all intermediate nodes execute only \texttt{ForwardException}, which appends diagnostic context and relays the exception upward without attempting local resolution.

\begin{algorithm}[htbp]
\caption{Bailout Protocol escalation algorithm.}
\label{alg:escalate}
\begin{algorithmic}[1]
\Procedure{Escalate}{$n$, trigger}
  \State $\text{ctx} \leftarrow$ \Call{BuildDiagnosticContext}{$n$, trigger}
  \State $\text{path} \leftarrow$ \Call{SelectPathway}{$n$}
  \State $\text{retry} \leftarrow 0$
  \State $T_{\text{prop}} \leftarrow T_{\text{prop}}^{\text{init}}$
  \State \Call{RecordToLedger}{$n$, ESCALATING, trigger, ctx}
  \While{$\text{retry} < N_{\max}$}
    \State $\text{ack\_received} \leftarrow$ \Call{SendToParent}{$n$, ctx, $T_{\text{prop}}$}
    \If{$\text{ack\_received} = \texttt{ACK}$}
      \State \Return \texttt{ACK\_RECEIVED}
    \EndIf
    \State $\text{path} \leftarrow$ \Call{SelectPathway}{$n$} \Comment{re-evaluate on retry}
    \State $T_{\text{prop}} \leftarrow \min(2 \cdot T_{\text{prop}},\; T_{\text{cap}})$
    \State $\text{retry} \leftarrow \text{retry} + 1$
    \State \Call{RecordToLedger}{$n$, RETRY, retry, path}
  \EndWhile
  \State \Comment{$N_{\max}$ retries exhausted — fire parent\_unreachable directly to $h(n)$}
  \State \Call{FireParentUnreachable}{$n$, ctx}
  \State \Return \texttt{UNREACHABLE\_ESCALATED}
\EndProcedure
\Statex
\Procedure{ForwardException}{$n$, ctx$,$ parent}
  \State $\text{ctx}' \leftarrow$ \Call{AppendDiagnosticContext}{$n$, ctx}
  \State \Call{RecordToLedger}{$n$, FORWARD, ctx'}
  \State $\text{ack\_received} \leftarrow$ \Call{SendToParent}{$n$, ctx$'$, $T_{\text{prop}}$}
  \If{$\neg\text{ack\_received}$}
    \State \Call{RecordToLedger}{$n$, PARENT\_NACK, ctx$'$}
  \EndIf
  \State \Return $\text{ack\_received}$
\EndProcedure
\end{algorithmic}
\end{algorithm}

\begin{algorithm}[htbp]
\caption{Subroutines supporting the escalation algorithm.}
\label{alg:subroutines}
\begin{algorithmic}[1]
\Function{BuildDiagnosticContext}{$n$, trigger}
  \State $\text{ctx} \leftarrow \emptyset$
  \State $\text{ctx.node} \leftarrow n$
  \State $\text{ctx.trigger} \leftarrow \text{trigger}$
  \State $\text{ctx.bounds\_state} \leftarrow$ \Call{GetBoundsState}{$n$}
  \State $\text{ctx.resolution\_attempts} \leftarrow$ \Call{GetResolutionAttempts}{$n$}
  \State $\text{ctx.hop\_count} \leftarrow 0$
  \State $\text{ctx.chain} \leftarrow [\,]$ \Comment{empty list — originating node records no predecessors}
  \State \Return ctx
\EndFunction
\Statex
\Function{AppendDiagnosticContext}{$n$, ctx}
  \State $\text{ctx}'.node \leftarrow n$
  \State $\text{ctx}'.bounds\_state \leftarrow$ \Call{GetBoundsState}{$n$}
  \State $\text{ctx}'.hop\_count \leftarrow \text{ctx.hop\_count} + 1$
  \State $\text{ctx}'.chain \leftarrow \text{ctx.chain} \;\concat\; [n]$ \Comment{append current node to chain}
  \State \Return $\text{ctx}'$
\EndFunction
\Statex
\Function{SelectPathway}{$n$}
  \State \Comment{Returns the lowest-latency functional pathway from the Pathway Health Registry}
  \State $\text{candidates} \leftarrow \bigl\{\,p \in \text{PHR.pathways} \mid \text{p.status} = \texttt{FUNCTIONAL}\,\bigr\}$
  \If{$\text{candidates} = \emptyset$}
    \State \Return $\text{PHR.default\_emergency\_pathway}$
  \EndIf
  \State \Return $\argmin_{p \in \text{candidates}} \; \text{PHR.latency\_estimate}(p)$
\EndFunction
\Statex
\Procedure{FireParentUnreachable}{$n$, ctx}
  \State \Comment{All normal pathways exhausted — alert $h(n)$ via emergency pathway}
  \State $\text{alert} \leftarrow$ \Call{BuildDiagnosticContext}{$n$, \texttt{parent\_unreachable}}
  \State $\text{alert.severity} \leftarrow 1.0$ \Comment{maximum severity — human must act immediately}
  \State $\text{alert.hop\_count} \leftarrow \text{ctx.hop\_count}$
  \State $\text{alert.chain} \leftarrow \text{ctx.chain}$
  \State \State \Call{RecordToLedger}{$n$, PARENT\_UNREACHABLE, alert}
  \State \Call{SendToHuman}{$h(n)$, alert, $T_{\text{emergency}}$}
\EndProcedure
\end{algorithmic}
\end{algorithm}

The \texttt{BuildDiagnosticContext} function (Algorithm~\ref{alg:subroutines}, line 1) assembles the complete diagnostic record that travels up the parent chain. It captures the node identifier, the raw trigger event, the Bounds Engine's internal state at the moment of bailout initiation, all resolution attempts made during the BOUNDED window, and the current hop count. The \texttt{hop\_count} field starts at 0 on the originating node and is incremented by each intermediate node via \texttt{AppendDiagnosticContext} (line 11), providing $h(n)$ with a direct measure of chain depth. The \texttt{chain} list accumulates node identifiers in order, enabling $h(n)$ to reconstruct the complete propagation path retrospectively from the Forensic Ledger.

The \texttt{SelectPathway} function (line 19) queries the Pathway Health Registry (PHR), a runtime data structure maintained by the Bailout Monitor that tracks the availability, latency, and error rate of each provisioned communication pathway to the human anchor. Pathways that fail $K$ consecutive health checks (default $K = 3$) are marked \texttt{DEGRADED} and excluded from selection. The PHR is updated by a background health-ping thread operating at scheduling priority equal to the Bailout Monitor, ensuring that pathway status remains current even under heavy agent workload. The \texttt{default\_emergency\_pathway} is a provisioned fallback — typically an out-of-band channel (e.g., direct SMS or email via a registered integration) — that is used when all normal pathways are degraded or unreachable.

\subparagraph{Pathway Health Registry.} The PHR is defined as a tuple:
\begin{equation}
\text{PHR} \;=\; \bigl(\text{pathways},\; \text{latency\_estimate},\; \text{error\_count},\; \text{health\_check\_timestamp}\bigr)
\end{equation}
where \texttt{pathways} is the set of all provisioned communication channels to $h(n)$. Before each \texttt{SendToParent} call, \texttt{SelectPathway} queries the PHR and selects the pathway with the lowest estimated latency among those with \texttt{status} = \texttt{FUNCTIONAL}. The latency estimate is maintained by the background health-ping thread using an exponentially weighted moving average of observed round-trip times.

% --------------------------------------------------------------
\subsection{Parent-Chain Bypass Mechanism}\label{sec:bypass}

The parent-chain bypass is the structural property that ensures the Bailout Protocol's escalation pathway contains no machine actors as terminal or intermediate absorption points between the originating node and $h(n)$. Formally:
\begin{equation}
\forall n \in \text{nodes},\; \forall m \in \text{chain}\bigl(n, h(n)\bigr) \setminus \{h(n)\},\; m \notin \text{machine\_actor} \;\lor\; \text{terminal}(m) = \text{false}. \tag{BYPASS}
\end{equation}
This is restated more directly: for every node $n$, the set of all intermediate nodes on the chain from $n$ to its human anchor $h(n)$ contains exactly zero machine actors. The only non-machine actor on the chain is $h(n)$ itself. This set contains exactly one element: the human anchor $h(n)$. No machine actor may appear in this set.

\subparagraph{Why bypass is architecturally necessary.} The bypass is not an optimization — it is a structural necessity that follows from the upstream accountability guarantee of the \bxthree{} Framework. Any machine actor $m$ that could absorb an exception in ESCALATING state would become a terminal node for that exception, and the accountability chain would terminate at $m$ rather than at a human. If $m$ then fails, or if $m$'s resolution is incorrect or insufficient, no human principal is present in the chain to detect or correct the failure. The system compounds failure instead of surfacing it. This is the precise failure mode the Bailout Protocol is designed to prevent.

Conventional escalation hierarchies — including tiered oversight models that permit absorption at intermediate tiers — make human oversight a high-tier option rather than a compulsory terminus. The Bailout Protocol forbids this absorption by architectural constraint: $h(n)$ is not one option among many — it is the \emph{only} permissible terminus for every unresolved exception. The parent-chain bypass skips intermediate nodes not solely for efficiency but because their presence in the terminus would break the accountability chain by introducing machine actors as potential terminal nodes.

\subparagraph{Enforcement mechanisms.} The bypass is enforced by two cooperating mechanisms that operate at different layers of the \bxthree{} architecture:

\begin{enumerate}
\item[\textbf{Fact Layer gate (enforcement at the authorization layer).}] The Fact Layer is the only mechanism capable of authorizing state transitions in the \bxthree{} Framework. When a node $n$ is in BOUNDED state and produces a resolution $a$, the resolution is written to the authorization ledger only after Fact Layer approval. The Bailout Monitor's pre-commit hook rejects any ledger entry for a resolution where the underlying trigger satisfies TC and $T_{\text{bounds}}$ has expired. This prevents the Bounds Engine from issuing a retroactive resolution after timeout to suppress warranted escalation. The Fact Layer gate operates identically on all machine actors: there is no privileged machine actor with the ability to bypass this check.

\item[\textbf{Immutable parent-pointer and human-root architecture (enforcement at the data model layer).}] The parent pointer $\alpha(n)$ and the human-root identifier $h(n)$ are stored in Fact Layer-managed memory and are read-only at the machine-actor level. A machine actor cannot redirect the parent chain, cannot insert itself as an intermediate resolver, and cannot modify the $humanRootId$ field that terminates the chain. Any attempted modification requires \purpose-layer authorization with a corresponding ledger entry — which cannot be completed without human approval if the modification would affect the bailout terminus. This architectural immutability ensures that the propagation path is fixed at provisioning and cannot be altered by any machine actor at runtime.
\end{enumerate}

Together, these two mechanisms ensure that BYPASS holds for all $n \in \text{nodes}$ under all execution conditions, including partial node failure, network partition, and deliberate attempt by a compromised machine actor to redirect the bailout pathway.

% --------------------------------------------------------------
\subsection{Timeout Handling}\label{sec:timeout-handling}

Timeouts in the Bailout Protocol serve a liveness function: they prevent the protocol from permanently stalling when a parent node is unreachable, a communication pathway is unresponsive, or the human anchor is temporarily unavailable. Each timeout has a defined scope, a defined scope-exit action, and an escalation path that prevents the protocol from deadlocking.

\subparagraph{Timeout definitions.}

\begin{enumerate}
\item[\textbf{$T_{\text{bounds}}$ — Bounds Engine resolution timeout.}] The maximum time a node may remain in BOUNDED state attempting autonomous resolution. Default: 60 seconds, provisioned per node class according to the expected latency of the node's resolution procedure. If $T_{\text{bounds}}$ expires without a Fact-Layer-approved resolution, the node transitions to BAIL\_PENDING $\rightarrow$ ESCALATING via T3 and T4. The Fact Layer cannot extend $T_{\text{bounds}}$ unilaterally; any extension requires human authorization via a separate ledger entry recorded before $T_{\text{bounds}}$ expires.

\item[\textbf{$T_{\text{prop}}$ — Per-hop propagation timeout.}] The maximum time node $n$ waits for an acknowledgment from its parent $p$ after sending an exception via \texttt{SendToParent}. Default initial value: 5 seconds. If $T_{\text{prop}}$ expires without ACK, the algorithm retries with exponential backoff: $T_{\text{prop}} \leftarrow \min(2 \cdot T_{\text{prop}},\; T_{\text{cap}})$. This backoff continues until either an ACK is received or $N_{\max}$ retries have been exhausted at $T_{\text{cap}}$.

\item[\textbf{$T_{\text{cap}}$ — Maximum propagation timeout value.}] The ceiling value for the exponential backoff of $T_{\text{prop}}$. Default: 60 seconds. After $N_{\max}$ consecutive retries at $T_{\text{cap}}$ without ACK, the algorithm fires a \texttt{parent\_unreachable} timeout event (T7), which is treated as a new trigger event propagating to $h(n)$ via the next available functional pathway tracked by the Pathway Health Registry.

\item[\textbf{$T_{\text{human}}$ — Human acknowledgment timeout.}] The maximum time the human anchor may remain in HUMAN\_ACKNOWLEDGED state without issuing a directive. Default: 300 seconds (5 minutes), configurable per deployment based on the expected human response time for the node's operational domain. If $T_{\text{human}}$ expires, the protocol issues a reminder notification to $h(n)$ and transitions the originating node to RESOLVED with an \texttt{unresolved} tag. The node becomes quiescent: no autonomous retry occurs, no further actions are attempted, and the \texttt{unresolved} tag in the Forensic Ledger signals that human re-engagement is required to clear the state. The protocol does not attempt to re-escalate automatically; automatic re-escalation would create a potential denial-of-service condition against the human anchor.

\item[\textbf{$T_{\text{health}}$ — Pathway health-check interval.}] The interval at which the background health-ping thread sends heartbeat messages along each provisioned pathway. Default: 30 seconds. Pathways that fail $K$ consecutive health checks are marked \texttt{DEGRADED} and excluded from \texttt{SelectPathway} until a health check succeeds. This prevents the algorithm from repeatedly selecting a known-failed pathway during active escalation.
\end{enumerate}

\begin{table}[htbp]
\caption{Timeout parameters, their start conditions, and their functional scopes.}
\label{tab:timeouts}
\begin{tabular}{@{}lcl@{}}
\toprule
\textbf{Timeout} & \textbf{Started at} & \textbf{Scope} \\
\midrule
$T_{\text{bounds}}$   & T1 fires (TC)                           & Bounds Engine autonomous resolution window      \\
$T_{\text{prop}}$     & T4 fires (first \texttt{SendToParent})   & Per-hop parent acknowledgment                    \\
$T_{\text{cap}}$      & T6 fires (retry)                        & Ceiling for $T_{\text{prop}}$ backoff           \\
$T_{\text{human}}$    & ESCALATING + $h(n)$ contact confirmed   & Human directive issuance window                 \\
$T_{\text{health}}$   & Continuous background thread            & Pathway availability monitoring                  \\
\bottomrule
\end{tabular}
\end{table}

\subparagraph{Liveness guarantee.} Together, $T_{\text{bounds}}$, $T_{\text{prop}}$, $T_{\text{cap}}$, and $T_{\text{human}}$ define a finite upper bound on the wall-clock time from trigger firing (T1) to human acknowledgment under any combination of parent unavailability and pathway degradation. This bound is a deployment parameter — not a protocol constant — derived from the network topology, the number of hops to the human anchor, and the provisioned $T_{\max}$ parameter:
\begin{equation}
T_{\max} \;=\; T_{\text{bounds}} \;+\; N_{\max} \cdot T_{\text{cap}} \;+\; T_{\text{human}}. \tag{T-BOUND}
\end{equation}
A deployment that cannot guarantee human acknowledgment within the required $T_{\max}$ under all failure scenarios must provision additional or higher-bandwidth pathways, or must reduce $T_{\text{prop}}$ and $T_{\text{cap}}$ to bring the worst-case bound within acceptable limits. The liveness guarantee is therefore not merely a property of the protocol specification but also a constraint on deployment topology: the Bailout Protocol guarantees that escalation completes within a bounded time \emph{given} a correctly provisioned deployment.
